% Hướng dẫn cho Người chuẩn bị bài (NCB). Dịch: pdflatex slides.tex (hai lần), sau scripts/test-all.sh
\documentclass[aspectratio=169,12pt]{beamer}
\usepackage{../pi-hd}
\newcommand\vaitro{Hướng dẫn cho Người chuẩn bị bài}
\begin{document}

\bia{Người chuẩn bị bài}{Biên tập đề và lời giải, ghi mục cần kiểm tra và xung đột, vẽ hình, thêm bài mới.}

\chuanbi

\chu{Nguyên tắc biên tập}{%
\begin{buoc}
\item \textbf{Không sửa lén}: mỗi sửa đổi nội dung toán là một dòng ghi — vị trí, trước, sau, lý do; chờ tác giả xác nhận.
\item \textbf{Bản gốc của tác giả giữ nguyên}; thầy cô chỉ sửa bản biên tập.
\item \textbf{Văn bản sạch}: đề và lời giải chỉ có \LaTeX; ghi chú để ở mục Cần kiểm tra, Xung đột.
\item \textbf{Nghi ngờ thì ghi, không tự sửa}: thành mục cần kiểm tra; khi cần, hỏi tác giả qua Văn phòng hoặc Phụ trách chuyên mục.
\end{buoc}
\medskip Mọi lần sửa đều có lịch sử (ai, khi nào, khác chỗ nào); hai người sửa cùng lúc không ghi đè lên nhau.}

\manhinh{dang-nhap}{Vào hệ thống}
\manhinh{danh-sach}{Danh sách bài}
\manhinh{loc}{Lọc và sắp xếp}
\manhinh{bai-moi-a}{Trang bài}
\manhinh[Sửa nhỏ: chính tả, định dạng. Sửa nội dung toán: công thức, đáp số, lập luận.]{sua-de}{Sửa đề, xem trước ngay}
\manhinh{hai-nguoi-sua}{Khi người khác vừa lưu}
\manhinh{nguon-chinh-sua}{Nguồn và chỉnh sửa}
\manhinh{hinh}{Hình vẽ TikZ}
\manhinh{dung-tu-dong}{Dựng hình tự động}
\manhinh{hinh-tren-trang}{Hình trên trang bài}
\manhinh{them-bai}{Thêm bài: tác giả và tệp}
\manhinh{them-bai-2}{Thêm bài: ảnh và lưu}
\manhinh{da-them}{Bài vừa thêm}

\chu{Khi có sự cố}{%
\begin{itemize}\setlength\itemsep{2mm}
\item \textbf{„Bài vừa được người khác sửa"}: bản của thầy cô vẫn còn trong ô; so với bản mới nhất bên dưới, gộp phần cần giữ, bấm Lưu lần nữa.
\item \textbf{„Cảnh báo hiển thị"} dưới đề: công thức có lỗi \LaTeX\ (thiếu ngoặc, lệnh lạ) — sửa rồi lưu.
\item \textbf{Hình „chưa dựng"} sau vài phút: trang Hình → Dựng ngay; hình lỗi hiện kèm thông báo của \LaTeX.
\item \textbf{Huỷ khi đang sửa}: bấm Huỷ hai lần (lần đầu nhắc còn thay đổi chưa lưu).
\end{itemize}}

\gopy

\chu{Tóm tắt}{%
\begin{buoc}
\item Lọc danh sách, mở bài cần xử lý.
\item Sửa nhỏ, hoặc sửa nội dung toán kèm vị trí và lý do.
\item Ghi mục cần kiểm tra, xung đột; đóng mục phải ghi kết quả.
\item Vẽ hình bằng TikZ — hệ thống tự dựng.
\item Thêm bài mới ở trang Thêm bài.
\end{buoc}}
\end{document}
